\documentclass[10pt,a4paper]{amsart}
\usepackage[margin=19mm]{geometry}
\usepackage{amsmath,amssymb,mathtools}
\usepackage[scr=rsfs,cal=euler]{mathalfa}
\usepackage{xfrac}
\usepackage[unicode,hidelinks,bookmarks=false]{hyperref}
\newcommand{\ie}{i.e.\ }
\newcommand{\cf}{cf.\ }
\newcommand{\resp}{respectively\ }
\newcommand{\Subsection}{\subsection}
\providecommand{\nobreakdash}{\nobreak\hbox{-}}
\setlength{\emergencystretch}{2em}
\multlinegap0pt
\usepackage{bm}
\usepackage{bbm}
\usepackage{dsfont}
\usepackage{xspace}
\usepackage{cancel}
\usepackage{mathtools}
\usepackage{amsthm}
\makeatletter
\@ifundefined{theorem}{\newtheorem{theorem}{Theorem}}{}
\@ifundefined{lemma}{\newtheorem{lemma}[theorem]{Lemma}}{}
\makeatother
\DeclareMathOperator{\C}{C}
\DeclareMathOperator{\Z}{Z}
\newcommand\XX{\mathfrak{X}}
\title[On a class of involutive Yang-Baxter groups]{On a class of involutive Yang-Baxter groups}
\author{H. Meng, A. Ballester-Bolinches, P. Jim\'enez-Seral}
\date{Source: arXiv:2509.06521v2 (CC BY 4.0)}
\begin{document}
\maketitle

\noindent\emph{Source and licence.} On a class of involutive Yang-Baxter groups. Version arXiv:2509.06521v2 (CC BY 4.0), distributed under the Creative Commons Attribution 4.0 International licence, as recorded in the arXiv licence metadata for this exact version and shown on the abstract page https://arxiv.org/abs/2509.06521v2. Full rights notice: licenses/arxiv-2509.06521-CC-BY-4.0.txt.\par\smallskip

\noindent\textit{Editorial scope.} Counted unit: the Lemma whose source label is \texttt{lem-elements} in the article (source0.tex lines 836--853, source label \texttt{lem-elements}), delivered together with the article's own complete original proof of it (source0.tex lines 854--948, one \texttt{proof} environment of 95 lines). No mathematics is written, repaired or re-derived: statement and proof are the article's own text line for line. Required context retained uncounted: none. Every definition, display and label of those lines is retained unchanged. Omitted from this package and not redistributed: 1--835, 949--1108. Nothing inside a retained excerpt is omitted or altered: every retained body is delivered exactly as the article's lines are written, so no omission receipt of any kind accompanies this package. Citations: the retained excerpts carry no citation command at all, so no bibliography entry of the article is transcribed and no reference is redistributed. Reference adaptation: none was needed and none was made; every reference inside the delivered unit resolves inside it, so no unresolved reference or citation is produced. Printed slips kept literal: no printed word, symbol, hypothesis or number of the counted statement or of its proof was altered. Layout and class: the article's own class is \texttt{article} with packages including fontenc, babel, inputenc, amsmath, amssymb, amsthm, hyperref; the wrapper is the lane's pinned \texttt{amsart} editorial wrapper at 10pt on A4, which restates the theorem-like environments the retained text needs and adds the article's own macro definitions that the retained text needs (\texttt{\textbackslash C}, \texttt{\textbackslash XX}, \texttt{\textbackslash Z}), with extra wrapper packages (bm, bbm, dsfont, xspace, cancel, mathtools). The delivered numbering is the unit's own. Assets: no retained span contains a figure environment, a graphics-inclusion command, a PDF-page inclusion, a table or a TikZ picture; no image file is copied into this package, the package declares no asset dependency and no image file is redistributed with it. Licence: version arXiv:2509.06521v2 of this article is distributed under the Creative Commons Attribution 4.0 International licence, as recorded in the arXiv licence metadata for this exact version and shown on the abstract page https://arxiv.org/abs/2509.06521v2. A full rights notice is delivered at licenses/arxiv-2509.06521-CC-BY-4.0.txt. Characters: every retained excerpt is pure ASCII, so the delivered engine, XeLaTeX with its default Latin Modern text font, reports no missing character.
\begin{lemma}\label{lem-elements}
Let $G=A_1A_2\cdots A_k$ be the product of the subgroups $A_i$,  $1 \leq i \leq k$, such that $A_i$ is normalised by $A_j$ for $1\leq i<j \leq k$. Suppose that
$$A_1A_2 \cdots A_i \cap A_{i+1}\cdots A_k \leq \C_G(A_iA_{i+1}),~i=1,2,\cdots,k-1.$$
Let $a_1a_2\cdots a_k=b_1b_2\cdots b_k$, where $a_i,b_i \in A_i$ for each $i$. Then
\begin{itemize}
\item[\emph{(1)}] For each $1\leq i \leq k-1$, we have that
$$\emph{\text{(1-1)}}~(b_1^{-1}a_1)(b_2^{-1}a_2)\cdots (b_i^{-1}a_i)=(b_{i+1}\cdots b_k)(a_{i+1}\cdots a_k)^{-1} \in \C_G(A_iA_{i+1});$$
$$\emph{\text{(1-2)}}~(a_{i+1}b_{i+1}^{-1})(a_{i+2}b_{i+2}^{-1})\cdots (a_kb_k^{-1})=(a_{1}\cdots a_{i})^{-1} (b_{1}\cdots b_{i}) \in \C_G(A_iA_{i+1});$$
\item[\emph{(2)}] $a_ib_i^{-1}, b_i^{-1}a_i \in \Z(A_i)$ for each $1 \leq i \leq k$; In particular if $b_1=b_2=\cdots=b_k=1$, then $a_i \in \Z(A_i).$
\item[\emph{(3)}] Let $h_i=a_{i+1}a_{i+2} \cdots a_{k}~(1\leq i \leq k-1)$ and $h_k=1$. Suppose that
$x_1x_2 \cdots x_k=1$ for some $x_i \in A_i, 1\leq i \leq k$. Then
$$({}^{h_1}x_1)\cdot ({}^{h_2}x_2) \cdots ({}^{h_k}x_k)=1. $$
\item[\emph{(4)}] Write $$\XX_j=\{(a_1,a_2 \cdots, a_j) \mid a_1a_2\cdots a_j=1, a_i \in A_i, \forall 1 \leq i \leq j\}$$
    for each $j=1,2,\cdots,k$. Then
    $$|A_1A_2\cdots A_j| \cdot |\XX_j|=|A_1|\cdot |A_2| \cdots |A_j|$$
for each $j=1,2,\cdots,k.$
\end{itemize}
\end{lemma}

\begin{proof}
(1) We will first prove $\text{(1-1)}$ by induction on $i$. For $i=1$,
$$b_1^{-1}a_1=(b_{2}\cdots b_k)(a_{2}\cdots a_k)^{-1} \in A_1 \cap A_{2}A_3\cdots A_k \leq \C_G(A_1A_{2}).$$
Now we assume that $\text{(1-1)}$ holds for $i-1$, that is,  $$(b_1^{-1}a_1)(b_2^{-1}a_2)\cdots (b_{i-1}^{-1}a_{i-1})=(b_{i}\cdots b_k)(a_{i}\cdots a_k)^{-1} \in \C_G(A_{i-1}A_{i}),$$
which implies that
$$b_i^{-1}(b_1^{-1}a_1)(b_2^{-1}a_2)\cdots (b_{i-1}^{-1}a_{i-1})a_{i+1}=(b_{i+1}\cdots b_k)(a_{i+1}\cdots a_k)^{-1}. $$
Note that $(b_1^{-1}a_1)(b_2^{-1}a_2)\cdots (b_{i-1}^{-1}a_{i-1}) \in \C_G(A_i)$ by induction, hence it centralizes $b_i^{-1} \in A_i$. As $A_{i+1}A_{i+2} \cdots A_k \leq G$,  we get that
$$(b_1^{-1}a_1)(b_2^{-1}a_2)\cdots (b_i^{-1}a_i)=(b_{i+1}\cdots b_k)(a_{i+1}\cdots a_k)^{-1}$$
belongs to $A_1A_2\cdots A_i \cap A_{i+1}\cdots A_k \leq \C_G(A_iA_{i+1})$, as desired.
We next show (1-2) by backward induction on $i$. For $i=k-1$, by hypothesis,
$$a_kb_k^{-1}=(a_{1}\cdots a_{k-1})^{-1} (b_{1}\cdots b_{k-1}) \in A_1A_2\cdots A_{k-1} \cap A_k \leq \C_G(A_{k-1}A_k).$$
We assume that it is true for $i+1$, that is,
$$(a_{i+2}b_{i+2}^{-1})(a_{i+3}b_{i+3}^{-1})\cdots (a_kb_k^{-1})=(a_{1}\cdots a_{i+1})^{-1} (b_{1}\cdots b_{i+1}) \in \C_G(A_{i+1}A_{i+2}),$$
which implies that
$$a_{i+1}(a_{i+2}b_{i+2}^{-1})(a_{i+3}b_{i+3}^{-1})\cdots (a_kb_k^{-1})b_{i+1}^{-1}=(a_{1}\cdots a_{i})^{-1} (b_{1}\cdots b_{i}).$$
Since $(a_{i+2}b_{i+2}^{-1})(a_{i+3}b_{i+3}^{-1})\cdots (a_kb_k^{-1}) \in \C_G(A_{i+1})$ by induction, which commutes with $b_{i+1}^{-1} \in A_{i+1}$, it follows that
$$(a_{i+1}b_{i+1}^{-1})(a_{i+2}b_{i+2}^{-1})\cdots (a_kb_k^{-1})=(a_{1}\cdots a_{i})^{-1} (b_{1}\cdots b_{i}),$$
which is in $A_1 A_2\cdots A_i \cap A_{i+1} \cdots A_k \leq
\C_G(A_iA_{i+1})$. Therefore Statement~$(1)$ is proved.
\\
\\
(2) Since that Equation (1-1) holds for $i=1,k-1$, we have that
$$b_1^{-1}a_1=(b_{2}\cdots b_k)(a_{2}\cdots a_k)^{-1} \in \C_G(A_1A_{2}) \cap A_1 \leq \Z(A_1);~\text{and}$$
$$(b_1^{-1}a_1)(b_2^{-1}a_2)\cdots (b_{k-1}^{-1}a_{k-1})=b_ka_k^{-1} \in \C_G(A_{k-1}A_{k}) \cap A_k \leq \Z(A_k).$$
Hence $b_i^{-1}a_i \in \Z(A_i)$ for $i=1$ and $k$. Moreover, for each $2\leq i \leq k-1$, it follows from Statement~(1) that
$$(b_1^{-1}a_1)(b_2^{-1}a_2)\cdots(b_{i-1}^{-1}a_{i-1}) (b_i^{-1}a_i) \in \C_G(A_iA_{i+1}) \leq \C_G(A_i);$$
and
$$(b_1^{-1}a_1)(b_2^{-1}a_2)\cdots(b_{i-1}^{-1}a_{i-1}) \in \C_G(A_{i-1}A_i) \leq \C_G(A_i).$$
Thus $b_i^{-1}a_i \in \C_G(A_i) \cap A_i=\Z(A_i)$ for $2\leq i \leq k-1$.  Hence we deduce that $b_i^{-1}a_i \in \Z(A_i)$ for each $1 \leq i \leq k$, and so $b_i^{-1}a_i=b^{-1}_i(a_ib_i^{-1})b_i \in \Z(A_i)$, as desired.
\\
\\
(3) In fact, we will prove
$$({}^{h_i}x_i)\cdot ({}^{h_{i+1}}x_{i+1}) \cdots ({}^{h_k}x_k)={}^{h_i}(x_ix_{i+1}\cdots x_k)$$
by backward induction on $i$. It is trivial for $i=k$ and we assume that
$$({}^{h_{i+1}}x_{i+1}) \cdots ({}^{h_k}x_k)={}^{h_{i+1}}(x_{i+1}\cdots x_k)$$
holds, where $i \leq k-1$.
Note that $h_i=a_{i+1}h_{i+1}$ and $$h_i^{-1}h_{i+1}=h_{i+1}^{-1}a_{i+1}^{-1}h_{i+1} \in A_{i+1}$$
since $h_{i+1}=a_{i+2}a_{i+3} \cdots a_{k} \in A_{i+2}A_{i+3}\cdots A_k$ normalises $A_{i+1}$ by hypothesis.
As $x_1x_2\cdots x_k=1$ and $A_1A_2\cdots A_i \leq G$ by hypothesis, we have
$x_{i+1}x_{i+2}\cdots x_k=(x_1x_2\cdots x_i)^{-1} \in A_1A_2\cdots A_i \cap A_{i+1} \cdots A_k \leq \C_G(A_iA_{i+1}).$
Hence $x_{i+1}x_{i+2}\cdots x_k$ commutes with $h_i^{-1}h_{i+1}$. Hence
$$
\begin{aligned}
({}^{h_{i}}x_{i})({}^{h_{i+1}}x_{i+1}) \cdots ({}^{h_k}x_k)&=({}^{h_{i}}x_{i}) \cdot {}^{h_{i+1}}(x_{i+1}x_{i+2}\cdots x_k)\\
&=h_ix_i(h_i^{-1}h_{i+1})(x_{i+1}x_{i+2}\cdots x_k)h_{i+1}^{-1}\\
&=h_ix_i(x_{i+1}x_{i+2}\cdots x_k)(h_i^{-1}h_{i+1})h_{i+1}^{-1}\\
&={}^{h_i}(x_ix_{i+1}x_{i+2}\cdots x_k),
\end{aligned}
$$
as desired. Consequently, $({}^{h_1}x_1)\cdot ({}^{h_{2}}x_{2}) \cdots ({}^{h_k}x_k)={}^{h_1}(x_1x_{2}\cdots x_k)=1$ and Statement~(3) holds.
\\
\\
(4) We prove Statement~(4) by induction on $j$. Clearly, we may assume that $j \geq 2$. Note that $A_1A_2 \cdots A_{j-1}$ is a subgroup of $G$. By induction we have
$$
\begin{aligned}
|A_1|\cdot |A_2| \cdots |A_{j-1}| \cdot |A_j|&=|A_1A_2\cdots A_{j-1}| \cdot |\XX_{j-1}| \cdot |A_j|\\
&= |A_1A_2\cdots A_{j-1}|\cdot |A_j| \cdot |\XX_{j-1}| \\
&=|A_1A_2\cdots A_j| \cdot |A_1A_2\cdots A_{j-1} \cap A_j| \cdot |\XX_{j-1}|
\end{aligned}
$$
Write $B=A_1A_2\cdots A_{j-1} \cap A_j$. We only have to show that
$|\XX_j|=|B| \cdot |\XX_{j-1}|$.
For $b \in A_j$, we consider the set
$$\XX_j(b)=\{(a_1,a_2,\cdots,a_{j-1},b) \mid a_1a_2\cdots a_{j-1}b=1, a_i \in A_i, \forall 1 \leq i \leq j-1\},$$
which is a subset of $\XX_j$.
 It is easy to see that $\XX_j(b) \neq \varnothing$ if and only if $b \in B$. Note that $\XX_j$ is the disjoint union of all $\XX_j(b)$ for $b \in B$. It implies that $|\XX_j|=\sum_{b \in B} |\XX_j(b)|.$

We claim that $|\XX_j(b)|=|\XX_{j-1}|$ for each $b \in B$. Fix an element $b \in B$, as $B=A_1A_2\cdots A_{j-1} \cap A_j$, we may assume that $b^{-1}=b_1b_2\cdots b_{j-1}$ for some $b_i \in A_i, 1\leq i \leq j-1$. Let
$\phi: \XX_j(b) \rightarrow \XX_{j-1}$ be a correspondence defined as follows:
$$\phi((a_1,a_2,\cdots, a_{j-1}, b))=(a_1b_1^{-1}, a_2b_2^{-1}, \cdots, a_{j-1}b_{j-1}^{-1}).$$
Since $a_1a_2\cdots a_{j-1}=b^{-1}=b_1b_2\cdots b_{j-1}$, by Statement~(1), it follows that $(a_1b_1^{-1})(a_2b_2^{-1})\cdots(a_{j-1}b_{j-1}^{-1})=1$. Hence

$$\phi((a_1,a_2,\cdots, a_{j-1}, b))=(a_1b_1^{-1}, a_2b_2^{-1}, \cdots, a_{j-1}b_{j-1}^{-1}) \in \XX_{j-1}.$$
Hence $\phi$ is actually a map. For each two elements $(a_1,a_2,\cdots,a_{j-1},b)$ and $(a_1',a_2',\cdots,a_{j-1}',b)$ in $\XX_j(b)$ with
$$\phi((a_1,a_2,\cdots,a_{j-1},b))=\phi((a_1',a_2',\cdots,a_{j-1}',b)),$$
we have that $a_ib_i^{-1}=a_i'b_i^{-1}$ for $1 \leq i \leq j-1$, and so $a_i=a_i'$. Hence $\phi$ is injective.  For each $(a_1,a_2,\cdots, a_{j-1}) \in \XX_{j-1}$,
we have $a_1a_2\cdots a_{j-1}=1$.
Write $l_0=1$, $l_i=a_1 a_2\cdots a_{i}, 1\leq i \leq j-1$ and note that $l_{j-1}=a_1 a_2\cdots a_{j-1}=1$. Then
$$l_i=a_1a_2\cdots a_{i}=(a_{i+1}\cdots a_{j-1})^{-1} \in A_1A_2\cdots A_i \cap A_{i+1} \cdots A_k$$
centralizes $b_i \in A_i$ by hypothesis.
As $a_{i}=l_{i-1}^{-1}l_{i}$ for $1\leq i \leq j-1$, we have that
$$
\begin{aligned}
&(a_1b_1)(a_2b_2)\cdots (a_{j-2}b_{j-2})(a_{j-1}b_{j-1})b\\
=&(l_1b_1)(l_1^{-1}l_2b_2)(l_2^{-1}l_3b_3) \cdots (l_{j-3}^{-1}l_{j-2}b_{j-2})(l_{j-2}^{-1}l_{j-1}b_{j-1})b\\
=&({}^{l_1}b_1)({}^{l_2}b_2)({}^{l_3}b_3) \cdots ({}^{l_{j-2}}b_{j-2}) b_{j-1}b\\
=&(b_1b_2\cdots b_{j-2}b_{j-1})b=b^{-1}b=1,
\end{aligned}
$$
which implies that $(a_1b_1,a_2b_2,\cdots,a_{j-1}b_{j-1}, b) \in \XX_j(b)$. Then
$$\phi((a_1b_1,a_2b_2,\cdots, a_{j-1}b_{j-1}, b))=(a_1, a_2, \cdots, a_{j-1}).$$
Hence $\phi$ is bijective and $|\XX_j(b)|=|\XX_{j-1}|$ for each $b \in B$. Consequently,
$$|\XX_j|=\sum_{b \in B} |\XX_j(b)|=|B||\XX_{j-1}|,$$
as desired. The proof of the lemma is now complete.
\end{proof}

\end{document}
